\subsection{子模的交与和}

对每个族 \((\mathbf{M}_{\iota})_{\iota \in I}\) （A-模 \(\mathbf{E}\) 的子模的族），交 \(\bigcap_{\iota \in I} \mathbf{E}_{\iota}\) 是 \(\mathbf{E}\) 的一个子模。若对每个 \(\iota \in I\) ， \(\phi_{\iota}\) 表示典范同态 \(\mathbf{E} \to \mathbf{E} / \mathbf{M}_{\iota}\) ，则 \(\bigcap_{\iota \in I} \mathbf{M}_{\iota}\) 是同态 \(\phi: x \mapsto (\phi_{\iota}(x))\) （ \(\mathbf{E}\) 到 \(\prod_{\iota \in I} (\mathbf{E} / \mathbf{M}_{\iota})\) 中的同态）的核，换言之，存在正合列

\[
0 \longrightarrow \bigcap_ {i \in I} M _ {i} \longrightarrow E \stackrel {\phi} {\longrightarrow} \prod_ {i \in I} (E / M _ {i}).\tag{27}
\]

线性映射 \(\phi\) 以及映射

\[
\mathrm{E} / \left(\bigcap_ {i \in I} \mathrm{M} _ {i}\right)\rightarrow \prod_ {i \in I} (\mathrm{E} / \mathrm{M} _ {i})
\]

（它是通过取商得到的）都称为典范映射。

特别地：

命题 8. 若子模族 \((\mathbf{M}_i)_{i\in I}\) （ \(\mathbf{E}\) 的子模族）的交约化为 0 ，则 \(\mathbf{E}\) 典范同构于 \(\prod_{i\in I}(\mathbf{E}/\mathbf{M}_i)\) 的一个子模。

给定 A-模 E 的一个子集 X，包含 X 的 E 的诸子模的交 F 称为由 X 生成的子模，而 X 称为 F 的生成集（或生成系）（I，§4，第 3 号）；对 E 的元素的一个族 \((a_{i})_{i\in I}\) ，由诸 \(a_{i}\) 组成的集合所生成的子模称为由族 \((a_{i})\) 生成的子模。

一个 A-模称为有限生成的，如果它有一个有限的生成集。

命题 9. 由 A-模 E 的元素族 \((a_{i})_{i\in I}\) 生成的子模是族 \((a_{i})\) 的线性组合的集合。

E 的包含所有 \(a_{i}\) 的每个子模也包含诸 \(a_{i}\) 的线性组合。反之，第 1 号的公式 (1) 和 (2) 证明诸 \(a_{i}\) 的线性组合的集合是 E 的一个子模，它显然包含所有 \(a_{i}\) ，因而是包含它们的最小子模。

推论 1. 设 \(u: \mathbf{E} \to \mathbf{F}\) 是线性映射， \(\mathbf{S}\) 是 \(\mathbf{E}\) 的一个子集，而 \(\mathbf{M}\) 是 \(\mathbf{E}\) 中由 \(\mathbf{S}\) 生成的子模。则 \(u(\mathbf{M})\) 是 \(\mathbf{F}\) 中由 \(u(\mathbf{S})\) 生成的子模。

特别地，E 的任何有限生成子模在 u 下的像是 F 的一个有限生成子模。

注记。若 \(u(x) = 0\) （对所有 \(x \in S\) ），则也有 \(u(x) = 0\) （对所有 \(x \in M\) ）。我们有时把这一结果称为“线性恒等式的延拓原理”或“线性延拓原理”。

特别地，要验证线性映射 \(u:E\to F\) 具有形式 \(v\circ\phi\) ，其中 \(v:E/M\to F\) 是线性的而 \(\phi:E\to E/M\) 是典范投影，只需验证 \(u(S)=0\) 。

推论 2. 设 \((\mathbf{M}_{t})_{t\in I}\) 是模 \(\mathbf{E}\) 的一个子模族；由该族的并生成的子模与形如 \(\sum_{i\in I}x_i\) 的和的集合恒同，其中 \((x_{i})_{i\in I}\) 取遍 \(\mathbf{E}\) 中元素的具有有限支集且满足 \(x_{i}\in \mathbf{M}_{i}\) （对所有 \(i\in I\) ）的族。

显然， \(\bigcup_{i\in I}M_i\) 中元素的每个线性组合都具有上述形式：逆命题是显然的。

由 E 的子模族 \((\mathbf{M}_{\iota})_{\iota\in I}\) 的并生成的 E 的子模称为该族 \((\mathbf{M}_{\iota})\) 的和，记作 \(\sum_{\iota\in I}\mathbf{M}_{\iota}\) 。若对所有 \(\iota\in I\) ， \(h_{\iota}\) 是典范单射 \(M_{\iota}\to E\) ，而 \(h:(x_{\iota})\mapsto\sum_{\iota}h_{\iota}(x_{\iota})\) 是 \(\bigoplus_{\iota\in I}M_{\iota}\) 到 E 中的、对应于族 \((h_{\iota})\) 的线性映射（第 6 号，命题 6），则 \(\sum_{\iota\in I}M_{\iota}\) 是 h 的像；换言之，存在正合列

\[
\bigoplus_ {i \in I} M _ {i} \xrightarrow {h} E \longrightarrow E / \left(\sum_ {i \in I} M _ {i}\right) \longrightarrow 0.\tag{28}
\]

推论 3. 若 \((\mathbf{M}_{\lambda})_{\lambda \in L}\) 是 A-模 \(\mathbf{E}\) 的子模的一个右定向族，则和 \(\sum_{\lambda \in L} \mathbf{M}_{\lambda}\) 与并 \(\bigcup_{\lambda \in L} \mathbf{M}_{\lambda}\) 恒同。

\(\bigcup_{\lambda \in L} M_{\lambda} \subset \sum_{\lambda \in L} M_{\lambda}\) 无需假设总成立；另一方面，对每个有限子族 \((M_{\lambda})_{\lambda \in J}\) （族 \((M_{\lambda})_{\lambda \in L}\) 的），根据假设存在 \(\mu \in L\) 使得 \(M_{\lambda} \subset M_{\mu}\) （对所有 \(\lambda \in J\) ），从而 \(\sum_{\lambda \in L} M_{\lambda} \subset M_{\mu}\) ，于是由推论 2 得 \(\sum_{\lambda \in L} M_{\lambda} \subset \bigcup_{\lambda \in L} M_{\lambda}\) 。

推论 4. 设 \(0 \to \mathbf{E} \xrightarrow{f} \mathbf{F} \xrightarrow{g} \mathbf{G} \to 0\) 是 A-模的一个正合列，S 是 E 的生成系，T 是 G 的生成系。若 \(\mathbf{T}'\) 是 F 的一个子集，使得 \(g(\mathbf{T}') = \mathbf{T}, \mathbf{T}' \cup f(\mathbf{S})\) 是 F 的一个生成系。

子模 \(\mathbf{F}'\) （ \(\mathbf{F}\) 中由 \(\mathbf{T}' \cup f(\mathbf{S})\) 生成的子模）包含 \(f(\mathbf{E})\) ，并且由于 \(g(\mathbf{F}')\) 包含 \(\mathbf{T}, g(\mathbf{F}') = \mathbf{G}\) ；因此 \(\mathbf{F}' = \mathbf{F}\) 。

推论 5. 在 A-模的正合列 \(0 \to \mathbf{E} \to \mathbf{F} \to \mathbf{G} \to 0\) 中，若 E 和 G 都是有限生成的，则 F 也是有限生成的。

命题 10. 设 M, N 是 A-模 E 的两个子模。则存在两个正合列

\begin{enumerate}
\def\labelenumi{(\arabic{enumi})}
\setcounter{enumi}{28}
\tightlist
\item
\end{enumerate}

\[
0 \longrightarrow \mathrm{M} \cap \mathrm{N} \stackrel {u} {\longrightarrow} \mathrm{M} \oplus \mathrm{N} \stackrel {i - j} {\longrightarrow} \mathrm{M} + \mathrm{N} \longrightarrow 0\tag{30}
\]

\[
0 \longrightarrow \mathrm{E} / (\mathrm{M} \cap \mathrm{N}) \stackrel {v} {\longrightarrow} (\mathrm{E} / \mathrm{M}) \oplus (\mathrm{E} / \mathrm{N}) \stackrel {p - q} {\longrightarrow} \mathrm{E} / (\mathrm{M} + \mathrm{N}) \longrightarrow 0
\]

其中 \(i:\mathbf{M}\to \mathbf{M} + \mathbf{N},j:\mathbf{N}\rightarrow \mathbf{M} + \mathbf{N}\) 是典范单射，

\[
p: \mathrm{E} / \mathrm{M} \rightarrow \mathrm{E} / (\mathrm{M} + \mathrm{N}) \quad and \quad q: \mathrm{E} / \mathrm{N} \rightarrow \mathrm{E} / (\mathrm{M} + \mathrm{N})
\]

是典范满射，而同态 u 和 v 定义如下：

若 \(f: \mathbf{M} \cap \mathbf{N} \to \mathbf{M} \to \mathbf{M} \oplus \mathbf{N}\) 和 \(g: \mathbf{M} \cap \mathbf{N} \to \mathbf{N} \to \mathbf{M} \oplus \mathbf{N}\) 是典范单射，则 \(u = f + g\) ，而若 \(r: \mathbf{E}/(\mathbf{M} \cap \mathbf{N}) \to \mathbf{E}/\mathbf{M} \to (\mathbf{E}/\mathbf{M}) \oplus (\mathbf{E}/\mathbf{N})\) 和

\[
s: \mathbf {E} / (\mathbf {M} \cap \mathbf {N}) \rightarrow \mathbf {E} / \mathbf {N} \rightarrow (\mathbf {E} / \mathbf {M}) \oplus (\mathbf {E} / \mathbf {N})
\]

是典范映射，则 \(v = r + s\) 。

我们证明 (29) 的正合性：显然 \(i-j\) 是满射且 \(u\) 是单射。另一方面，说 \((i-j)(x,y)=0\) （其中 \(x\in\mathbf{M}\) ， \(y\in\mathbf{N}\) ），就是说 \(i(x)-j(y)=0\) ，从而 \(i(x)=j(y)=z\in\mathbf{M}\cap\mathbf{N}\) ，于是按定义有 \(x=f(z)\) ， \(y=g(z)\) ，这就证明了 \(\operatorname{Ker}(i-j)=\operatorname{Im}u\) 。

我们证明 (30) 的正合性：显然 \(p - q\) 是满射。另一方面，说 \(v(t) = 0\) （对 \(t \in \mathbf{E} / (\mathbf{M} \cap \mathbf{N})\) ），就是说 \(r(t) = s(t) = 0\) ，从而 \(t\) 是模 ( \(\mathbf{M} \cap \mathbf{N}\) ) 的、某个元素 \(z \in \mathbf{E}\) 的类，该元素模 \(\mathbf{M}\) 和模 \(\mathbf{N}\) 的类均为零，这蕴含 \(z \in \mathbf{M} \cap \mathbf{N}\) ，故 \(t = 0\) 。最后，说 \((p - q)(x, y) = 0\) （其中 \(x \in \mathbf{E} / \mathbf{M}\) ， \(y \in \mathbf{E} / \mathbf{N}\) ），就是说 \(p(x) = q(y)\) ，或者说存在两个元素 \(z'\) , \(z''\) （属于 \(\mathbf{E}\) ），它们模 \(\mathbf{M}\) 和模 \(\mathbf{N}\) 的类分别是 \(x\) 和 \(y\) ，并且满足 \(z' - z'' \in \mathbf{M} + \mathbf{N}\) 。于是存在 \(t' \in \mathbf{M}\) , \(t'' \in \mathbf{N}\) 使得 \(z' - z'' = t' - t''\) ，由此

\[
z ^ {\prime} - t ^ {\prime} = z ^ {\prime \prime} - t ^ {\prime \prime} = z.
\]

设 \(w\) 为 \(z; r(w)\) 模 (M ∩ N) 的类；它是 \(z\) 模 M 的类，从而也是 \(z'\) 的类，即 \(x\) ；类似地 \(s(w) = y\) ，这就完成了 \(\operatorname{Ker}(p - q) = \operatorname{Im} v\) 的证明。

